\section*{The Propers: Themes and Movement}

\begin{massbox}{A king whose strength is humble}
The officially correlated strand presents one royal paradox: God's reign arrives in a king whose strength is humble and whose yoke gives rest. Beside it, Romans continues its own apostolic course, locating Christian freedom in the Spirit who raises bodies and puts sinful deeds to death, while the offering and Communion texts carry these themes toward purification, salvation, and enduring praise.
\end{massbox}

\begin{studybox}{Mercy, meekness, Spirit, praise}
\textbf{1. Gathered by mercy --- Entrance; Collect.} Temple mercy and the Son's humility gather and raise a fallen people.\par\smallskip
\textbf{2. The humble King speaks --- First Reading; Psalm; Acclamation; Gospel.} Peaceful rule and filial revelation form the officially correlated strand, and each text keeps its own history.\par\smallskip
\textbf{3. Freedom under the Spirit --- Second Reading.} Paul's semi-continuous strand calls hearers to Spirit-enabled mortification and bodily hope.\par\smallskip
\textbf{4. Purified, fed, praising --- Prayer over Offerings; Communion; Prayer after Communion.} The appointed prayers and antiphons order disciples toward salvation and praise.
\end{studybox}

\subsection*{1. Mercy gathers those who cannot raise themselves\enspace\elementref{Entrance; Collect}}

The Entrance begins with received mercy rather than religious achievement. Its Temple setting makes the gathering ecclesial, while praise reaches beyond one sanctuary. The Collect then attributes the world's raising to the Son's humility. These texts establish dependence before the readings specify kingship, Spirit, and discipleship; because they belong to Week XIV in every cycle, the relation is liturgical juxtaposition.

\subsection*{2. Royal power takes the form of meekness\enspace\elementref{First Reading; Psalm; Acclamation; Gospel}}

Zechariah removes the normal machinery of conquest from the king's arrival: humble riding, abolished weapons, peace proclaimed outward. Psalm 145 answers by praising the Lord whose kingdom is manifested through compassion, fidelity, and support. Matthew then gives the correlation its Christological center. Jesus does not merely recommend gentleness; as Son he reveals the Father, and as the meek teacher he summons disciples beneath his yoke.

\subsection*{3. Rest and mortification are not rivals\enspace\elementref{Second Reading; Gospel}}

Romans remains a semi-continuous strand, yet it prevents sentimental readings of rest. The Spirit's indwelling generates obligation: evil deeds are to be put to death. Chrysostom and Augustine insist that this mortification concerns disordered action rather than hatred of the created body. The Gospel's light yoke likewise names obedient discipleship, not inactivity. Their liturgical co-occurrence permits this comparison.

\subsection*{4. Gift becomes purified life and enduring praise\enspace\elementref{Prayer over Offerings; Communion; Prayer after Communion}}

The Prayer over the Offerings asks that worship purify and redirect conduct. Communion Antiphon A stresses tasting and refuge; alternative B repeats Christ's invitation to the burdened. The final prayer asks that great gifts issue in salvation and unceasing praise, making reception a beginning of transformed life rather than a terminus.
